\documentclass[10pt,a4paper]{amsart}
\usepackage[margin=19mm]{geometry}
\usepackage{amsmath,amssymb,mathtools}
\usepackage[scr=rsfs,cal=euler]{mathalfa}
\usepackage{xfrac}
\usepackage[unicode,hidelinks,bookmarks=false]{hyperref}
\newcommand{\ie}{i.e.\ }
\newcommand{\cf}{cf.\ }
\newcommand{\resp}{respectively\ }
\newcommand{\Subsection}{\subsection}
\providecommand{\nobreakdash}{\nobreak\hbox{-}}
\setlength{\emergencystretch}{2em}
\multlinegap0pt
\usepackage{bm}
\usepackage{bbm}
\usepackage{dsfont}
\usepackage{xspace}
\usepackage{cancel}
\usepackage{mathtools}
\usepackage{amsthm}
\makeatletter
\@ifundefined{Theorem}{\newtheorem{Theorem}{Theorem}}{}
\@ifundefined{Proposition}{\newtheorem{Proposition}[Theorem]{Proposition}}{}
\makeatother
\title[The asymptotic charges of Curtright dual graviton and Curtri]{The asymptotic charges of Curtright dual graviton and Curtright extensions of BMS algebra}
\author{Federico Manzoni}
\date{Source: arXiv:2602.20037v2 (CC BY 4.0)}
\begin{document}
\maketitle

\noindent\emph{Source and licence.} The asymptotic charges of Curtright dual graviton and Curtright extensions of BMS algebra. Version arXiv:2602.20037v2 (CC BY 4.0), distributed under the Creative Commons Attribution 4.0 International licence, as recorded in the arXiv licence metadata for this exact version and shown on the abstract page https://arxiv.org/abs/2602.20037v2. Full rights notice: licenses/arxiv-2602.20037-CC-BY-4.0.txt.\par\smallskip

\noindent\textit{Editorial scope.} Counted unit: the Proposition whose source label is \texttt{prop:ellipticity-Delta-tilde} in the article (source0.tex lines 2325--2331, source label \texttt{prop:ellipticity-Delta-tilde}), delivered together with the article's own complete original proof of it (source0.tex lines 2333--2403, one \texttt{proof} environment of 71 lines). No mathematics is written, repaired or re-derived: statement and proof are the article's own text line for line. Required context retained uncounted: lap+ (lines 2310--2313); lap1 (lines 2319--2322). Every definition, display and label of those lines is retained unchanged. Omitted from this package and not redistributed: 1--2309, 2314--2318, 2323--2324, 2332--2332, 2404--2798. Nothing inside a retained excerpt is omitted or altered: every retained body is delivered exactly as the article's lines are written, so no omission receipt of any kind accompanies this package. Citations: the retained excerpts carry no citation command at all, so no bibliography entry of the article is transcribed and no reference is redistributed. Reference adaptation: none was needed and none was made; every reference inside the delivered unit resolves inside it, so no unresolved reference or citation is produced. Printed slips kept literal: no printed word, symbol, hypothesis or number of the counted statement or of its proof was altered. Layout and class: the article's own class is \texttt{article} with packages including fontenc, yfonts, comment, jheppub, longtable, array, macro, appendix, amsmath, amssymb, tikz, pgfkeys, pgfopts, xcolor; the wrapper is the lane's pinned \texttt{amsart} editorial wrapper at 10pt on A4, which restates the theorem-like environments the retained text needs and adds the article's own macro definitions that the retained text needs (none), with extra wrapper packages (bm, bbm, dsfont, xspace, cancel, mathtools). The delivered numbering is the unit's own. Assets: no retained span contains a figure environment, a graphics-inclusion command, a PDF-page inclusion, a table or a TikZ picture; no image file is copied into this package, the package declares no asset dependency and no image file is redistributed with it. Licence: version arXiv:2602.20037v2 of this article is distributed under the Creative Commons Attribution 4.0 International licence, as recorded in the arXiv licence metadata for this exact version and shown on the abstract page https://arxiv.org/abs/2602.20037v2. A full rights notice is delivered at licenses/arxiv-2602.20037-CC-BY-4.0.txt. Characters: every retained excerpt is pure ASCII, so the delivered engine, XeLaTeX with its default Latin Modern text font, reports no missing character.
\begin{equation}
\Delta_{\delta^{(2)}}:=\delta^{(2)}\big(\delta^{(2)}\big)^{\dagger}+\big(\delta^{(2)}\big)^{\dagger}\delta^{(2)}
\label{lap+}
\end{equation}

\begin{equation}
\Delta:= d^{(1)}(d^{(1)})^{\dagger}+(d^{(1)})^{\dagger}d^{(1)}+d^{(2)}(d^{(2)})^{\dagger}+(d^{(2)})^{\dagger}d^{(2)}=\Delta_1 + \Delta_2\ .
\label{lap1}
\end{equation}

\begin{Proposition}\label{prop:ellipticity-Delta-tilde}
Let
\begin{equation*}\label{eq:Delta-tilde}
\widetilde{\Delta}\;:=\;\Delta^{2}+\Delta_{\delta^{(2)}}
\end{equation*}
the fourth-order operator where $\Delta$ and $\Delta_{\delta^{(2)}}$ are defined in \eqref{lap1} and \eqref{lap+}. Then $\widetilde{\Delta}$ is elliptic on $\Omega^{p\otimes q}(M)$ and so is on a Young projected subspace.
\end{Proposition}

\begin{proof}
Let $\xi\in T_x^*M$.
The principal symbols of the first-order operators are
\[
\sigma_1\!\bigl(d^{(i)}\bigr)(x,\xi)=\varepsilon_i(\xi):=\xi\wedge_{(i)},\qquad
\sigma_1\!\bigl((d^{(i)})^\dagger\bigr)(x,\xi)=\iota_i(\xi):=\iota^{(i)}_{\xi^\sharp}.
\]
It follows that the principal symbol of each $\Delta^{(i)}$ is
\begin{equation*}\label{eq:principal-symbol-slot-laplacian}
\sigma_2\!\bigl(\Delta^{(i)}\bigr)(x,\xi)=|\xi|^2\,\mathrm{Id},
\quad\Rightarrow\quad
\sigma_2(\Delta)(x,\xi)=2|\xi|^2\,\mathrm{Id},
\end{equation*}
and therefore
\begin{equation*}\label{eq:principal-symbol-Delta-square}
\sigma_4\!\bigl(\Delta^2\bigr)(x,\xi)=\bigl(2|\xi|^2\bigr)^2\,\mathrm{Id}=4|\xi|^4\,\mathrm{Id}.
\end{equation*}
Moreover,
\[
\sigma_2\!\bigl(\delta^{(2)}\bigr)(x,\xi)=\varepsilon_2(\xi)\varepsilon_1(\xi),
\qquad
\sigma_2\!\bigl((\delta^{(2)})^\dagger\bigr)(x,\xi)=\iota_1(\xi)\iota_2(\xi),
\]
so $\sigma_4(\Delta_{\delta^{(2)}})(x,\xi)$ is positive semidefinite but generally not invertible since
\[
\sigma_4\!\bigl(\delta^{(2)}(\delta^{(2)})^\dagger\bigr)(x,\xi)
=\varepsilon_2\varepsilon_1\iota_1\iota_2
=|\xi|^4\,\Pi^{(1)}_{+}(\xi)\,\Pi^{(2)}_{+}(\xi),
\]
\[
\sigma_4\!\bigl((\delta^{(2)})^\dagger\delta^{(2)}\bigr)(x,\xi)
=\iota_1\iota_2\varepsilon_2\varepsilon_1
=|\xi|^4\,\Pi^{(1)}_{-}(\xi)\,\Pi^{(2)}_{-}(\xi).
\]
where
\[
\Pi^{(i)}_{+}(\xi):=\frac{1}{|\xi|^2}\,\varepsilon_i(\xi)\,\iota_i(\xi),\qquad
\Pi^{(i)}_{-}(\xi):=\frac{1}{|\xi|^2}\,\iota_i(\xi)\,\varepsilon_i(\xi).
\]
Therefore
\begin{equation}\label{eq:symbol-Delta-delta2}
\sigma_4\!\bigl(\Delta_{\delta^{(2)}}\bigr)(x,\xi)
=|\xi|^4\Bigl(\Pi^{(1)}_{+}(\xi)\Pi^{(2)}_{+}(\xi)+\Pi^{(1)}_{-}(\xi)\Pi^{(2)}_{-}(\xi)\Bigr),
\end{equation}
which has kernel on $\Omega^{p_1 \otimes p_2}(M)$.
In particular, $\Delta_{\delta^{(2)}}$ is not elliptic by itself.

For $\xi\neq 0$, the principal symbol of $\widetilde{\Delta}$ is the sum
\[
\sigma_4(\widetilde{\Delta})(x,\xi)
=
4|\xi|^4\,\mathrm{Id}+\sigma_4(\Delta_{\delta^{(2)}})(x,\xi);
\]
moreover,
\begin{itemize}
    \item since $\Delta$ is a second-order elliptic operator, its symbol $\sigma_{\Delta}(x, \xi)$ is strictly positive definite. Consequently, the symbol of $\Delta^2$ is also strictly positive definite:
    \begin{equation}
    \langle \sigma_{\Delta^2}(x, \xi) v, v \rangle > 0, \quad \forall v \neq 0, \, \forall \xi \neq 0.
    \end{equation}
    
    \item the operator $\Delta_{\delta^{(2)}}$ is non-elliptic, meaning its symbol $\sigma_{\Delta_{\delta^{(2)}}}(x, \xi)$ may have a non-trivial kernel in certain directions. However, its algebraic structure (it is defined by the Hodge type structure) ensures it is semi-definite positive:
    \begin{equation}
    \langle \sigma_{\Delta_{\delta^{(2)}}}(x, \xi) v, v \rangle \geq 0, \quad \forall v \neq 0, \, \forall \xi \neq 0.
    \end{equation}
\end{itemize}
Therefore,
\begin{equation}
\langle \sigma_{\widetilde{\Delta}}(x, \xi) v, v \rangle = \langle \sigma_{\Delta^2}(x, \xi) v, v \rangle + \langle \sigma_{\Delta_{\delta^{(2)}}}(x, \xi) v, v \rangle > 0, \quad \forall v \neq 0, \, \forall \xi \neq 0;
\end{equation}
since $\sigma_{\widetilde{\Delta}}(x, \xi)$ is strictly positive definite, all its eigenvalues are strictly positive, ensuring it is invertible. This proves that $\widetilde{\Delta}$ is an elliptic operator.
\end{proof}

\end{document}
